% !TEX root = ../tfg.tex

Este capítulo explica los conceptos necesarios para entender cómo funciona el sistema y cómo se evalúa.
No pretende ser un manual, y por eso cada concepto se presenta solo hasta donde se utiliza en el trabajo y termina indicando dónde aparece.
Se parte de las ideas generales de las redes neuronales, después se presentan los tipos de capa que se emplean y por último se explican las herramientas de audio, de visión y de evaluación.

\section{Aprendizaje automático y redes neuronales}
\label{sec:fund-aprendizaje}

Programar un sistema que reconozca un golpe de pádel con reglas escritas a mano resulta casi imposible, porque no existe una descripción sencilla de lo que distingue una derecha de un revés en términos de coordenadas.
El aprendizaje automático ataca el problema desde el lado contrario.
En lugar de describir la solución se muestran al sistema muchos ejemplos con su respuesta correcta, y es él quien deduce qué patrones separan unos casos de otros~\cite{goodfellow2016deep}.
Todos los modelos de este trabajo aprenden de este modo.

Una red neuronal es un modelo formado por capas de unidades sencillas encadenadas.
Cada unidad calcula una suma ponderada de sus entradas y le aplica una función no lineal, de modo que la salida de una capa pasa a ser la entrada de la siguiente.
Para una capa completa, esto equivale a multiplicar la entrada por una matriz de pesos, sumarle un vector de sesgos y aplicar la función no lineal, que puede ser por ejemplo la que anula los valores negativos.

Los pesos y los sesgos de todas las capas son los parámetros del modelo, y son lo que se ajusta al entrenar.
Su número da una idea del tamaño de la red.
En este trabajo va desde unos \num{115000} parámetros en el detector de golpes hasta casi 1,8 millones en el clasificador y en la red de la pelota.

Entrenar consiste en encontrar los parámetros con los que el modelo se equivoca menos.
Para medir cuánto se equivoca se define una función de pérdida, que compara lo que el modelo predice con la respuesta correcta y devuelve un número mayor cuanto peor es la predicción.
Entrenar equivale entonces a buscar los parámetros que minimizan esa pérdida sobre los ejemplos disponibles.

El método habitual para minimizarla es el descenso del gradiente.
En cada paso se calcula en qué dirección crece más la pérdida respecto de cada parámetro, que es lo que indica su gradiente, y se mueven los parámetros ligeramente en la dirección contraria.
El tamaño de ese paso lo fija la tasa de aprendizaje.
Si es demasiado grande el entrenamiento oscila sin llegar a ningún resultado, y si es demasiado pequeña avanza con mucha lentitud.

Calcular el gradiente con todos los ejemplos a la vez sería costoso, por lo que cada paso se hace con un pequeño subconjunto elegido al azar, que se llama lote.
Una pasada completa por todos los ejemplos se denomina época, y el entrenamiento repite varias de ellas.
Los optimizadores modernos, además, ajustan por separado el tamaño del paso de cada parámetro.
AdamW~\cite{loshchilov2019adamw}, que es el que se utiliza en este trabajo, es uno de ellos y añade un pequeño decaimiento de los pesos para evitar que crezcan sin control.

La pérdida que se utiliza en la red de la pelota y en el clasificador es la entropía cruzada.
Esta función compara la probabilidad de la salida del modelo con la respuesta correcta, y con ello castiga más los errores cometidos con mucha seguridad.
Cuando los casos que interesan son muy escasos, como la pelota, que ocupa muy pocas celdas del mapa frente al fondo, la entropía cruzada queda dominada por los demás.
Para evitarlo se multiplica el error de cada celda por un peso mayor donde está la pelota, lo que da la expresión de la \cref{eq:perdida-pelota}.

Si hay varias clases, como en el clasificador del tipo de golpe, la salida es una probabilidad por clase, y esas probabilidades las produce la función softmax, que se explica en la \cref{sec:fund-atencion}.
En ese caso la pérdida vale $-\sum_c t_c \log p_c$.
Donde tc vale 1 para la clase correcta y 0 para las demás, de modo que el modelo es castigado por la probabilidad que da a ella.

Un modelo con muchos parámetros puede acabar memorizando los ejemplos con los que se entrena en lugar de aprender lo que tienen en común.
A este fenómeno se le llama sobreajuste, y hace que la pérdida baje mucho con los datos de entrenamiento y mucho menos con datos nuevos.
Por tanto, el rendimiento de un modelo nunca se mide con los ejemplos con los que se ha entrenado, sino con otros que se han reservado de antemano.

Además, el resultado de un entrenamiento depende de factores aleatorios, como los valores iniciales de los pesos o el orden de los lotes.
Dos entrenamientos del mismo modelo no dan exactamente lo mismo, y por eso en este trabajo el detector de golpes se entrena tres veces y se informa de la media (\cref{sec:resultados-instante}).

Estas ideas son comunes a todos los modelos del sistema.
Lo que cambia entre ellos es la forma de las capas.

\section{Redes para imágenes y secuencias}
\label{sec:fund-redes}

Las capas descritas hasta ahora conectan cada entrada con todas las unidades, y esto resulta poco eficiente cuando la entrada es una imagen.
Una imagen de 1080p tiene más de dos millones de píxeles, de modo que una capa que diera a cada unidad un peso por píxel sería inmanejable.

La convolución resuelve este problema.
Consiste en deslizar por la imagen un filtro pequeño y calcular en cada posición una suma ponderada de los valores que cubre.
Como el mismo filtro se aplica en todas las posiciones, sus pesos se comparten y el número de parámetros es pequeño.
Una capa usa muchos filtros, y cada uno produce su propio mapa de resultados, llamado canal.

Una red convolucional (\acs{cnn}) apila varias de estas capas.
Las primeras detectan patrones sencillos y las siguientes los combinan en otros más complejos, porque cada capa trabaja sobre lo que han detectado las anteriores.
Entre capas suele reducirse el tamaño de la imagen, lo que obliga a la red a quedarse con lo esencial y permite que cada unidad abarque cada vez una zona mayor.

En este trabajo las capas convolucionales aparecen en el detector de sonido, que las aplica sobre el espectrograma (\cref{sec:instante}) y en la red de la pelota (\cref{sec:pelota}).

Para detectar la pelota es necesario tener un valor para cada zona de la imagen y por lo tanto se emplea una red de arquitectura U-Net~\cite{ronneberger2015unet}.
Un codificador comprime la imagen paso a paso, de modo que la red pueda reconocer qué hay en ella, y un decodificador la amplía de nuevo hasta recuperar un tamaño parecido al original.
Al comprimir se pierde detalle sobre la posición exacta, que el decodificador recupera con conexiones directas a las capas del codificador de su mismo tamaño.
El resultado es un mapa del tamaño de la imagen, o de una fracción de él, que en este trabajo es el mapa de calor de la pelota (\cref{sec:fund-vision}).

Las secuencias, como el movimiento de un jugador, pueden tratarse con una convolución parecida pero en una sola dimensión.
El filtro se desliza a lo largo del tiempo, de modo que cada paso de la salida se calcula a partir de él y de sus vecinos.
Esta es la idea de las redes convolucionales temporales (\acs{tcn})~\cite{bai2018tcn}, que convierten posiciones en movimiento porque recogen cómo cambian los valores de un paso al siguiente.

Para que cada paso abarque más tiempo sin aumentar el número de pesos se utiliza la dilatación, que consiste en separar los puntos del filtro dejando huecos entre ellos.
Si una capa tiene un filtro de tamaño $k$ y una dilatación $d$, cada paso de su salida cubre $(k-1)d+1$ pasos de la entrada.
Con varias capas, la zona total que abarca un paso de la salida se llama campo receptivo y vale $1 + \sum (k-1)d$ de todas las capas.
La \acs{tcn} del clasificador tiene dos capas con $k = 5$ y dilataciones 1 y 3, de modo que cada paso abarca 17 pasos de la entrada.

Otra forma de tratar secuencias son las redes recurrentes (\acs{rnn}), que las leen paso a paso y van guardando un estado interno que resume lo visto hasta el momento.
En cada paso el estado se actualiza con la entrada nueva.
La unidad recurrente con puertas (\acs{gru})~\cite{cho2014gru} controla esa actualización con una puerta que decide cuánto del estado anterior se conserva y cuánto se sustituye por información nueva.

Una \acs{rnn} normal solo conoce el pasado de cada paso.
En una versión bidireccional se lee la secuencia dos veces, una hacia delante y otra hacia atrás, y se combinan los dos resultados.
Así cada paso se interpreta teniendo en cuenta lo que ocurre antes y después, algo útil para distinguir un golpe de un ruido que se le parece.

Estas dos ideas se combinan en las redes convolucionales recurrentes (\acs{crnn})~\cite{adavanne2017sed}.
Las capas convolucionales detectan en cada instante el patrón local de frecuencias y las recurrentes lo interpretan con ayuda del contexto temporal.
Es la arquitectura del detector de golpes por sonido (\cref{sec:instante}).

\section{Atención y transformers}
\label{sec:fund-atencion}

Las \acs{rnn} procesan la secuencia paso a paso, y eso dificulta que un paso lejano influya en el actual.
En un golpe esto es un inconveniente, porque lo que ocurre en el impacto depende de cómo se preparó el jugador un segundo antes.
El mecanismo de atención resuelve este problema porque permite que cada paso consulte directamente a todos los demás, sin importar la distancia que los separe~\cite{vaswani2017attention}.

La idea puede entenderse como una consulta a una base de datos.
Cada paso de la secuencia genera tres vectores, que se llaman consulta, clave y valor.
La consulta de un paso expresa qué información busca, la clave de otro paso describe qué información ofrece y el valor es la información en sí.
Cuanto más se parecen una consulta y una clave, más peso recibe el valor correspondiente.

Para toda la secuencia a la vez, las consultas, las claves y los valores se reúnen en las matrices $Q$, $K$ y $V$.
Se obtienen multiplicando la entrada por tres matrices de pesos, que la red aprende durante el entrenamiento.
La atención se calcula entonces como:
\begin{equation}
  \mathrm{Atenci\acute{o}n}(Q, K, V) = \mathrm{softmax}\left( \frac{Q K^{\top}}{\sqrt{d_k}} \right) V,
  \label{eq:atencion}
\end{equation}
donde $d_k$ es el tamaño de las claves.

El producto $Q K^{\top}$ mide el parecido entre cada consulta y cada clave.
La función softmax convierte esos parecidos en pesos que suman uno, y estos pesos se aplican a los valores.
Así cada paso de la salida es una mezcla de los valores de todos los pasos, con más peso en los que su consulta considera relevantes.

La función softmax transforma un conjunto de valores cualesquiera en probabilidades.
Para cada valor, calcula la exponencial de su puntuación y la divide entre la suma de las exponenciales de todas las clases.
La puntuación más alta se lleva la mayor parte.
En el clasificador, softmax convierte los valores del último bloque en la probabilidad de cada tipo de golpe (\cref{sec:salida-clasificador}).

Con una sola atención cada paso repartiría su interés de una única manera, pero un golpe contiene varias relaciones a la vez.
La atención con varias cabezas repite el cálculo anterior en paralelo, cada vez con sus propias matrices de pesos, de modo que cada cabeza puede especializarse en un tipo de relación.
Los resultados de todas ellas se concatenan y se combinan con una última matriz de pesos.
En el clasificador hay seis cabezas.

La atención tiene una limitación importante, y es que no sabe en qué orden están los pasos.
Esto es necesario para clasificar un golpe, y por tanto es esencial saber el orden de la preparación, el impacto y el cierre.
Por ello, a cada paso se le suma una codificación posicional, que le da una firma de su posición dentro de la secuencia.

Se utiliza la codificación sinusoidal del Transformer, en la que cada posición recibe un vector formado por senos y cosenos de distinta frecuencia.
Posiciones cercanas reciben vectores parecidos y posiciones lejanas vectores distintos, y la red puede aprovechar esa diferencia para razonar sobre el orden.

Una capa del transformer combina la atención con una pequeña red densa que se aplica a cada paso por separado.
Cada una de las dos partes va precedida de una normalización y su resultado se suma a su propia entrada, de modo que la capa solo tiene que aprender qué añadir a lo que ya recibe.
Apilando varias capas se obtiene un codificador, que en el clasificador tiene dos capas para la secuencia temporal y una para la interacción.

Para obtener un único vector que resuma toda la secuencia se añade al principio un token de clase.
Es un vector sin contenido propio que la red aprende como cualquier otro peso.
Como participa en la atención con todos los demás pasos, a la salida del codificador contiene un resumen de la secuencia, y es ese vector el que se utiliza para decidir.

La atención cruzada aplica la misma operación a dos secuencias distintas, la postura y la pelota.
Las consultas proceden de una de ellas y las claves y los valores de la otra, de modo que cada paso de la primera pregunta a la segunda qué momentos suyos son relevantes.
En el clasificador la postura consulta a la pelota, y eso resuelve la ambigüedad que tiene la postura sola (\cref{sec:red-clasificador}).
Las matrices se obtienen como indica la \cref{eq:cruzada}.

El último bloque del clasificador es un perceptrón multicapa (\acs{mlp}), que no es más que varias capas como las de la \cref{sec:fund-aprendizaje} apiladas una tras otra.
Recibe los resúmenes de las secuencias y devuelve una puntuación por cada clase, que la función softmax convierte en probabilidades.

\section{Audio y espectrogramas}
\label{sec:fund-audio}

El detector de golpes no trabaja con el audio tal como se graba, y para entender el motivo hay que empezar por cómo se representa el sonido.
Un micrófono registra la presión del aire en cada instante, y el resultado es una señal que oscila a lo largo del tiempo.
Para guardarla en un ordenador esa señal se mide muchas veces por segundo, y el número de medidas por segundo se llama frecuencia de muestreo.
Los vídeos de este trabajo tienen el audio a 48~kHz, es decir, \num{48000} medidas cada segundo.

Esta representación tiene un inconveniente, y es que mezcla todas las frecuencias en un único valor por instante.
El golpe de una pala y el murmullo del público ocupan la misma señal, y es difícil distinguirlos mirando solo cómo sube y baja.
Lo que los separa es qué frecuencias contiene cada uno, es decir, si el sonido es grave o agudo y cuánta energía aporta cada zona del rango.
Para ver esa información hay que pasar de la señal en el tiempo a su descomposición en frecuencias.

La herramienta matemática para hacerlo es la transformada de Fourier, que descompone una señal en ondas de distintas frecuencias y dice cuánta energía aporta cada una.
Solo tiene sentido hasta la mitad de la frecuencia de muestreo, que con 48~kHz son 24~kHz, porque no pueden representarse frecuencias mayores.

Aplicar la transformada a todo el audio de un punto diría qué frecuencias contiene, pero no cuándo aparece cada una.
Un golpe dura unas pocas milésimas de segundo, y su huella quedaría diluida entre un minuto de sonido.
Por eso se divide el audio en trozos cortos que se solapan entre sí y se calcula la transformada de cada trozo por separado.

Esta técnica se llama transformada de Fourier de tiempo corto.
Cada trozo se multiplica antes por una ventana que suaviza sus bordes, para que el corte brusco no introduzca frecuencias que no existen en el sonido original.
En este trabajo cada trozo tiene 2048 muestras, unos 43~ms, y entre un trozo y el siguiente hay un salto de 1024, de modo que cada trozo se solapa con el anterior en la mitad.

El espectrograma $S(m, k)$ es la energía de ese resultado para el trozo $m$ y la frecuencia $k$, es decir, el cuadrado de su módulo.
Se representa como una imagen en la que el eje horizontal es el tiempo, el vertical son las frecuencias y el brillo indica la energía.
Cada columna corresponde a un trozo y equivale a unos 21~ms de audio, que es el salto entre trozos, y cada fila a una de las 1025 frecuencias, separadas unos 23~Hz entre sí.

El espectrograma permite reconocer a simple vista distintos tipos de sonido.
Un sonido de frecuencia estable, como un silbido, aparece como una línea horizontal porque mantiene su energía constante durante un tiempo.
El impacto de una pala es lo contrario, porque dura muy poco pero reparte su energía por muchas frecuencias a la vez.
De ahí que un golpe se vea como una raya vertical, que es lo que debe aprender a reconocer el detector (\cref{fig:espectrogramas}).

El espectrograma tiene todavía un inconveniente práctico, y es que cada columna tiene 1025 valores, muchos más de los que la red necesita.
Además, el oído no percibe todas las frecuencias con la misma finura.
Distingue muy bien una diferencia de 100~Hz entre dos sonidos graves, pero apenas nota la misma diferencia entre dos agudos.
La escala Mel imita esta percepción, y convierte las frecuencias en una escala en la que los graves quedan más separados que los agudos.

Para reducir el número de valores se agrupan las frecuencias en bandas, mediante unos filtros triangulares repartidos a intervalos iguales en la escala Mel.
Como esa escala crece más despacio en las frecuencias altas, las bandas graves son estrechas y las agudas anchas, de modo que se conserva más detalle donde el oído lo distingue mejor.
La energía de cada banda $b$ es la suma de la energía de las frecuencias que cubre su filtro $H_b(k)$.

Además, las energías de un audio varían en un rango enorme, desde el murmullo de fondo hasta el golpe más fuerte, y el oído percibe la intensidad de forma aproximadamente logarítmica.
Por esta razón, se aplica el logaritmo a la energía de cada banda, con lo que se comprimen las diferencias y los golpes suaves dejan de quedar tapados por los fuertes:
\begin{equation}
  L(m, b) = \log\!\Bigl( \varepsilon + \sum_{k} H_b(k)\, S(m, k) \Bigr).
  \label{eq:logmel}
\end{equation}
El término $\varepsilon$ es un valor muy pequeño que evita calcular el logaritmo de cero.
El resultado es el espectrograma log-Mel.

Con esto se obtiene la entrada del detector.
En este trabajo se utilizan 40 bandas, de modo que cada columna del espectrograma tiene 40 valores y equivale a unos 21~ms de audio.
La red analiza secuencias de 256 columnas, que abarcan algo más de cinco segundos, y puede así ver el golpe junto con el contexto que lo rodea (\cref{sec:instante}).

\section{Visión por computador}
\label{sec:fund-vision}

Se presentan a continuación tres conceptos de imagen que utiliza el sistema.
El primero es una medida del solape entre dos cajas, el segundo es la manera de representar la posición de un objeto pequeño y el tercero es la transformación que permite pasar de la imagen a la pista real.
Se aplican en la \cref{sec:jugadores}, en la \cref{sec:pelota} y en la \cref{sec:asignacion}.

Para seguir a una persona a lo largo de un vídeo hay que decidir, en cada fotograma, cuál de las cajas detectadas corresponde a quien ya se estaba siguiendo.
Una caja es el rectángulo ajustado alrededor de un objeto y se describe con cuatro números.
Una forma natural de resolverlo es emparejar la caja de un fotograma con la del siguiente según cuánto se solapan, porque una persona se mueve poco entre dos fotogramas consecutivos y su caja casi no cambia de sitio.
Esa medida se llama intersección sobre unión y es el cociente entre el área que comparten dos cajas $A$ y $B$ y el área que cubren entre las dos:
\begin{equation}
  \mathrm{IoU}(A, B) = \frac{|A \cap B|}{|A \cup B|}.
  \label{eq:iou}
\end{equation}
Vale 1 cuando las cajas coinciden exactamente y 0 cuando no se tocan.
Cuanto más rápido se mueve una persona, menor es el valor, y por eso conviene fijar cuánto solape mínimo se exige para aceptar un emparejamiento.
El seguimiento completo, con todo lo que añade a esta idea, se explica en la \cref{sec:jugadores}.

La pelota plantea el problema contrario.
Es tan pequeña que ocupa unos pocos píxeles, y una caja a su alrededor aporta poco más que el punto donde está.
En estos casos resulta más cómodo trabajar con un mapa de calor, que es una cuadrícula con el tamaño de la imagen, o una fracción de ella, en la que cada celda contiene un valor entre 0 y 1 que indica cuánta probabilidad hay de que el objeto esté en ese lugar.

Para que una red aprenda a producir esos mapas hay que darle, junto a cada fotograma, el mapa correcto que debería devolver.
Marcar solo la celda de la pelota con un 1 y dejar el resto a 0 sería demasiado exigente, porque un acierto a una celda de distancia se trataría igual que uno a la otra punta de la imagen.
Por eso se coloca una gaussiana centrada en la posición $(x_0, y_0)$ de la pelota, que forma una mancha suave con su valor máximo en el centro y que decae poco a poco hacia los lados:
\begin{equation}
  t(x, y) = \exp\!\left( -\frac{(x - x_0)^2 + (y - y_0)^2}{2\sigma^2} \right).
  \label{eq:gaussiana}
\end{equation}
El parámetro $\sigma$ controla lo ancha que es la mancha.

El último concepto permite pasar de lo que se ve en la imagen a posiciones reales sobre la pista.
El suelo es un plano y, cuando se mira desde un punto de vista oblicuo, su imagen queda deformada por la perspectiva, de modo que los rectángulos se convierten en trapecios.
La homografía es la transformación que relaciona un plano con su imagen.
Se escribe con una matriz $H$ de $3 \times 3$ que actúa sobre coordenadas homogéneas, es decir, sobre un punto $(x, y)$ al que se añade un tercer valor igual a 1:
\begin{equation}
  \begin{pmatrix} x' \\ y' \\ w \end{pmatrix} = H \begin{pmatrix} x \\ y \\ 1 \end{pmatrix}, \qquad
  (X, Y) = \left( \frac{x'}{w},\, \frac{y'}{w} \right).
  \label{eq:homografia}
\end{equation}
El punto $(X, Y)$ es la posición correspondiente en el otro plano, que en este trabajo es la pista medida en metros.

Bastan cuatro puntos de la imagen con su posición real para determinar la matriz $H$, pero el sistema utiliza seis para que los pequeños errores al marcarlos pesen menos (\cref{fig:marcado-pista}).

Esta transformación solo es correcta para los puntos que están sobre el plano del suelo.
Un objeto en el aire se proyecta sobre la pista en un lugar que no es el suyo.
Una pelota que se encuentra a dos metros de altura, por ejemplo, parecería estar en otra parte de la pista, y por eso la homografía no se utiliza para situarla en metros.
Se aplica en cambio a los pies de los jugadores, que sí están en el suelo (\cref{sec:asignacion}).

\section{Evaluación}
\label{sec:fund-evaluacion}

Entrenar un modelo es solo la mitad del trabajo, porque queda saber cuánto se puede confiar en él.
Esta sección explica las medidas que se utilizan para comprobarlo y cómo se estima la incertidumbre de cada cifra.
Casi todos los resultados del \cref{cap:resultados} se apoyan en estos conceptos.

Cuando un modelo debe decidir si algo ocurre o no, como que haya un golpe en un instante, hay tres casos que importan.
Un verdadero positivo es un golpe real que el modelo detecta, y un falso positivo es una detección de un golpe que no existe.
Un falso negativo es un golpe real que el modelo no llega a detectar.
Estos casos son la base de todas las medidas que siguen.

La precisión indica qué parte de lo que el modelo señala es correcto, y el \textit{recall} qué parte de lo que existía ha conseguido encontrar.
Una precisión alta significa que casi todo lo que el modelo detecta es real, y un \textit{recall} alto que apenas se le escapa nada.

Ambas medidas se complementan, porque un modelo puede mejorar una a costa de la otra.
Detectar solo los casos más evidentes da una precisión muy alta pero deja escapar muchos golpes, y señalar cualquier cosa sospechosa encuentra casi todos los golpes pero llena el resultado de falsas alarmas.
Para resumirlas en un solo número se utiliza el F1, que es su media armónica:
\begin{equation}
  F_1 = \frac{2 \cdot \mathrm{Precisi\acute{o}n} \cdot \mathrm{Recall}}{\mathrm{Precisi\acute{o}n} + \mathrm{Recall}}
  \label{eq:f1}
\end{equation}
A diferencia de la media normal, la armónica queda cerca del valor más bajo, de modo que un F1 alto exige que las dos medidas lo sean.

Cuando lo que se mide son instantes en el tiempo hay que decidir cuándo una detección coincide con un golpe real, porque nunca lo hacen de forma exacta.
Se fija una tolerancia, que es de 250~ms en el detector de sonido, y se acepta una detección si cae dentro de ella.
Cada golpe real solo puede emparejarse con una detección, para que un mismo acierto no se cuente dos veces.
Con la pelota se hace lo mismo, pero la tolerancia es una distancia de 15 píxeles.

Cuando hay varias clases, como en el tipo de golpe, la medida más sencilla es el \textit{accuracy}, que es la proporción de golpes bien clasificados.
Esta medida tiene un inconveniente cuando las clases están desequilibradas, porque un modelo que acierte solo las más frecuentes obtiene una cifra alta aunque ignore las demás.
En este trabajo el saque tiene unas ocho veces menos ejemplos que la derecha, y por eso el \textit{accuracy} se acompaña del macro-F1.

El macro-F1 calcula el F1 de cada clase por separado y hace la media de todos ellos.
De este modo todas las clases pesan lo mismo, y una clase poco frecuente que se clasifica mal baja el resultado tanto como una muy frecuente.

Toda cifra se calcula con una muestra limitada de datos, y con otra muestra saldría algo distinta.
Por tanto, una cifra sin una medida de su incertidumbre dice poco, y la forma habitual de darla es un intervalo de confianza.
Un intervalo del 95~\% es un rango que contendría el valor verdadero en el 95~\% de las ocasiones si el estudio se repitiera con otras muestras, y cuanto más ancho es, más incierta resulta la cifra.

El método de Monte Carlo consiste en aproximar un resultado repitiendo muchas veces un experimento aleatorio, igual que se estimaría la probabilidad de que salga un seis lanzando un dado en lugar de calcularla.
El método \textit{bootstrap}~\cite{efron1979bootstrap} lo aplica al cálculo de ese intervalo sin suponer ninguna forma concreta para los datos, y el experimento que repite parte de una idea sencilla, que es utilizar los propios datos como si fueran la población completa.
Se forma una muestra nueva extrayendo al azar, con repetición, tantos elementos como tiene la original, y se recalcula la cifra sobre ella.
Al repetirlo muchas veces se obtiene una colección de valores que muestra cuánto varía la cifra según la muestra.

Si se realizan $B$ remuestreos y $\hat{\theta}^{*}_1, \ldots, \hat{\theta}^{*}_B$ son los valores obtenidos, ordenados de menor a mayor, el intervalo del 95~\% queda delimitado por los percentiles 2,5 y 97,5 de esa colección:
\begin{equation}
  \mathrm{IC}_{95\%} = \left[\, \hat{\theta}^{*}_{(0{,}025\,B)},\ \hat{\theta}^{*}_{(0{,}975\,B)} \,\right].
  \label{eq:bootstrap-ic}
\end{equation}

Calcular el resultado de todas las muestras posibles sería inviable, y por eso el \textit{bootstrap} genera \num{10000} remuestreos al azar, un número suficiente para que los percentiles apenas cambien si se repite el cálculo.

Al remuestrear hay que tener en cuenta que los golpes de un mismo punto no son independientes, porque comparten jugadores, pista y momento del partido.
Remuestrear golpes sueltos trataría cada uno como una observación nueva y daría intervalos demasiado estrechos, por lo que se remuestrean grupos completos, que son los puntos o los torneos según el caso.
